\section{Introduction}

A nanoparticle image is not an intrinsic measurement of ``modality quality''. It is the result of a scene, a configured instrument, a forward model, detector units, noise assumptions, and a derivative convention. Syniscopy makes those choices explicit by putting multiple configured microscopes on one shared scene parameterization and then emitting the Fisher-information diagnostics, status flags, and provenance needed to interpret comparisons under that parameterization.

We define a \emph{microscope} as the unit that the Cram\'{e}r--Rao bound is computed on: a chosen contrast modality together with that instrument's own complete parameter set. The modality selects the physics backend and contrast mechanism (for example a fluorescence emission backend, or an electron-transport backend); the instrument-level parameters --- numerical aperture, pixel pitch, pupil sampling, wavelength, detector noise, dose, and backend-specific settings --- belong to the microscope, not to the modality. Consequently two microscopes can select the same modality and still differ entirely in their configurations, and a parameter that is meaningless for one modality is simply absent from microscopes that do not use it. Comparing modalities directly is then a derived special case: hold one instrument fixed and vary only the contrast modality, with each variant evaluated using only the parameters its modality consumes.

The contribution is not a new microscope forward model for every modality. It is a shared-scene information contract: one latent particle/sample state is rendered through multiple declared microscope configurations, differentiated on the same physical parameter frame, and reported with sampling, convergence, validation, fusion, normalization, and supervision metadata. The paper-facing outputs are lateral and axial Cramer--Rao diagnostics, \(\mathrm{SE}(3)\) rank and symmetry checks, detected-quanta-normalized comparisons, algebraic fusion and serial-allocation diagnostics, and supervision masks whose positive/ignore labels are gated by measurement support. Lateral localization uses a step-free spectral derivative; axial and orientation diagnostics keep the explicit rerendered perturbation checks required by those state axes.

We use the term \emph{contract} for an evaluation view: Contract-LP fixes lateral particle-position diagnostics, Contract-LZ fixes axial diagnostics, Contract-Q fixes detected-quanta normalization, and Contract-NR fixes native-regime source-use auditing. A contract is not a benchmark leaderboard. It states what is being optimized and which artifacts must pass before a row can be treated as validated.

The manuscript is therefore organized around the minimum load-bearing evidence: the shared-scene formulation, the symmetry/fusion consequences, the matched-microscope-packet implementation, the supervision policy, and demonstrations whose generated artifacts expose sampling, convergence, and provenance status.
